\documentclass{beamer}
\mode<presentation>
\usepackage{amsmath}
\usepackage{amssymb}
%\usepackage{advdate}
\usepackage{graphicx}
\graphicspath{{./figs/}}
\usepackage{adjustbox}
\usepackage{subcaption}
\usepackage{enumitem}
\usepackage{multicol}
\usepackage{mathtools}
\usepackage{listings}
\usepackage{url}
\def\UrlBreaks{\do\/\do-}
\usetheme{Boadilla}
\usecolortheme{lily}
\setbeamertemplate{footline}
{
  \leavevmode%
  \hbox{%
  \begin{beamercolorbox}[wd=\paperwidth,ht=2.25ex,dp=1ex,right]{author in head/foot}%
    \insertframenumber{} / \inserttotalframenumber\hspace*{2ex} 
  \end{beamercolorbox}}%
  \vskip0pt%
}
\setbeamertemplate{navigation symbols}{}

\providecommand{\nCr}[2]{\,^{#1}C_{#2}} % nCr
\providecommand{\nPr}[2]{\,^{#1}P_{#2}} % nPr
\providecommand{\mbf}{\mathbf}
\providecommand{\pr}[1]{\ensuremath{\Pr\left(#1\right)}}
\providecommand{\qfunc}[1]{\ensuremath{Q\left(#1\right)}}
\providecommand{\sbrak}[1]{\ensuremath{{}\left[#1\right]}}
\providecommand{\lsbrak}[1]{\ensuremath{{}\left[#1\right.}}
\providecommand{\rsbrak}[1]{\ensuremath{{}\left.#1\right]}}
\providecommand{\brak}[1]{\ensuremath{\left(#1\right)}}
\providecommand{\lbrak}[1]{\ensuremath{\left(#1\right.}}
\providecommand{\rbrak}[1]{\ensuremath{\left.#1\right)}}
\providecommand{\cbrak}[1]{\ensuremath{\left\{#1\right\}}}
\providecommand{\lcbrak}[1]{\ensuremath{\left\{#1\right.}}
\providecommand{\rcbrak}[1]{\ensuremath{\left.#1\right\}}}
\theoremstyle{remark}
\newtheorem{rem}{Remark}
\newcommand{\sgn}{\mathop{\mathrm{sgn}}}
\providecommand{\abs}[1]{\left\vert#1\right\vert}
\providecommand{\res}[1]{\Res\displaylimits_{#1}} 
\providecommand{\norm}[1]{\lVert#1\rVert}
\providecommand{\mtx}[1]{\mathbf{#1}}
\providecommand{\mean}[1]{E\left[ #1 \right]}
\providecommand{\fourier}{\overset{\mathcal{F}}{ \rightleftharpoons}}
%\providecommand{\hilbert}{\overset{\mathcal{H}}{ \rightleftharpoons}}
\providecommand{\system}[1]{\overset{\mathcal{#1}}{ \longleftrightarrow}}
%\providecommand{\system}{\overset{\mathcal{H}}{ \longleftrightarrow}}
	%\newcommand{\solution}[2]{\textbf{Solution:}{#1}}
%\newcommand{\solution}{\noindent \textbf{Solution: }}
\providecommand{\dec}[2]{\ensuremath{\overset{#1}{\underset{#2}{\gtrless}}}}
\newcommand{\myvec}[1]{\ensuremath{\begin{pmatrix}#1\end{pmatrix}}}
\let\vec\mathbf

\lstset{
%language=C,
frame=single, 
breaklines=true,
columns=fullflexible
}

\numberwithin{equation}{section}

\title{2.10.11}
\author{ai25btech11022 - Narshitha}
% \maketitle
% \newpage
% \bigskip
\begin{document}
{\let\newpage\relax\maketitle}
\renewcommand{\thefigure}{\theenumi}
\renewcommand{\thetable}{\theenumi} 
\textbf{Question}:\\
 Construct a triangle ABC with side $\vec{BC}=6cm,\angle{B}=45,\angle{A}=105$ \\
 \textbf{Solution}:\\
 Let position vector of $\vec{B}$ be $\myvec{0\\0}$ and a,b,c be the sides opposite to the vertices A,B and C\\
 Given a=6cm;
 \begin{align}
     \vec{C}=\myvec{6\\0}\\
     \vec{A}=\myvec{c\cos{\angle{B}}\\c\sin{\angle{B}}}\\
     \vec{A}=\myvec{\frac{c}{\sqrt{2}}\\\frac{c}{\sqrt{2}}}\\
     \angle{A}+\angle{B}+\angle{C}=180\\
     \therefore \angle{C}=30
 \end{align}
 in $\Delta ABC$
 \begin{align}
     b\cos{\angle{C}} +c\cos{\angle{B}} =a\\
     b\sin{\angle{C}}-c\sin{\angle{B}} =0\\
     \myvec{\cos{\angle{C}} && \cos{\angle{B}}\\ \sin{\angle{C}} && -\sin{\angle{B}}}\myvec{b\\c}=\myvec{6\\0}
 \end{align}
 using augmented matrix 
 \begin{align}
     \myvec{\cos{\angle{C}} && \cos{\angle{B}} && \vrule && 6\\ \sin{\angle{C}} && -\sin{\angle{B}} && \vrule && 0}\\
     \myvec{\sqrt{3}/2 && 1/\sqrt{2}&& \vrule && 6 \\1/2 && -1/\sqrt{2} &&  \vrule && 0}
 \end{align}
 $R_2 \longrightarrow \sqrt{3}R_2 -R_1$
 \begin{align}
     \myvec{\sqrt{3}/2 && 1/\sqrt{2}&& \vrule && 6 \\0 && -\sqrt{3}-1/\sqrt{2} &&  \vrule && -6}\\
     \implies c= \frac{6\sqrt{2}}{\sqrt{3}+1}\\
   \therefore  \vec{A}=\myvec{6\brak{\sqrt{3}-1}\\6\brak{\sqrt{3}-1}}
 \end{align}
 \frametitle{Graphical Representation}
\begin{center}
    \includegraphics[width=0.7\linewidth]{figs/fig1.png}
\end{center}
\end{document}
